% Claim proof file for row claim_3_33 -- "AsymmetricSeparoidAxioms".
% Auto-generated stub by scaffold/solve_chapter.py at row start. A manager
% agent dedicated to writing the TeX proof fills in the proof body below.
%
% This file is a sibling of `main.tex` inside `<subsection>/tex/`. The
% `subfiles` package lets it render both standalone and as part of
% `main.tex` via `\subfile{...}`.
%
% The statement is restated above the proof so the file is self-contained
% when rendered on its own. The proof must:
%   - Stay close to the lecture notes' proof if one exists (always search
%     the LN first for a `\begin{proof}` block immediately following the
%     claim; copy / lightly edit when present).
%   - Otherwise use the LN paradigm and cite prior claims by their `ref`.
%   - Be self-contained enough that a mathematician can follow it without
%     re-reading the LN.
%   - Have no `sorry`, `omitted`, or "obvious" shortcuts.

\documentclass[main]{subfiles}

\begin{document}
% ref: claim_3_33 (proof, with statement restated for readability)
% title: AsymmetricSeparoidAxioms
\def\rowref{claim\_3\_33}
\phantomsection\label{claim_3_33}

% --- statement (restated) ---------------------------------------------------
\begin{DefThm}[(Asymmetric) separoid axioms for $i\sigma$-separation/$id$-separation]
    \label{d-sep-separoid-axioms}
    Let $G=(J,V,E,L)$ be a CDMG and $A,B,C,D \ins J \cup V$ subsets of nodes.
Then the ternary relations $\iPerp = \idPerp_G$ and $\iPerp = \isPerp_G$ satisfy the following rules:
\begin{enumerate}[label=\alph*)]
    \item Extended Left Redundancy: 
    \item[] $D \ins A \implies D \iPerp  B \given A$.
    \item $J$-Restricted Right Redundancy: 
    \item[] $A \iPerp  \emptyset \given C \cup J$ always holds.
    \item $J$-Inverted Right Decomposition: 
    \item[] $A \iPerp  B \given C \implies A \iPerp  J \cup B \given C$.
	\item Left Decomposition:  
	\item[]$A \cup D \iPerp  B \given C \implies D \iPerp  B \given C$.
  \item Right Decomposition:  	
    \item[] $A \iPerp  B \cup D\given C \implies A \iPerp  D \given C$.
    \item Left Weak Union:
    \item[] $A \cup D \iPerp  B  \given C \implies A \iPerp  B \given D \cup C$.
  \item Right Weak Union:\item[] $A \iPerp  B \cup D \given C \implies A \iPerp  B \given D \cup C$.
\item Left Contraction:\item[] 
 $ (A \iPerp  B \given D \cup C) \land (D \iPerp  B \given C) \implies A \cup D \iPerp  B \given C$.
\item Right Contraction:
\item[] $ (A \iPerp  B \given D \cup C) \land (A \iPerp  D \given C)  \implies A \iPerp  B \cup D \given C$.
\item Right Cross Contraction:
\item[] $ (A \iPerp  B \given D \cup C) \land (D \iPerp  A \given C) \implies A  \iPerp  B \cup D \given C$.
\item Flipped Left Cross Contraction: 
\item[] $ (A \iPerp  B \given D \cup C) \land (B \iPerp  D \given C) \implies B \iPerp  A \cup D \given C$.
\end{enumerate}
In particular, we have the equivalences:
$$ (A \iPerp  B \cup D \given C) \quad\iff\quad (A \iPerp  B \given D \cup C) \quad\land\quad (A \iPerp  D \given C)  ,$$
%and:
$$ (A \cup D \iPerp  B \given C) \quad\iff\quad (A \iPerp  B \given D \cup C) \quad\land\quad (D \iPerp  B \given C).$$
%Finally, if $J = \emptyset$ then 
We also get:
\begin{enumerate}[resume,label=\alph*)]
    \item $J$-Restricted Symmetry: 
	\item[] $A \iPerp  B \given C \cup J \implies B \iPerp A \given C \cup J$. 
\end{enumerate}
For the special case of $J=\emptyset$ we have thus (unrestricted) Symmetry.
\end{DefThm}

% --- proof ------------------------------------------------------------------
\begin{proof}
% TODO: write the proof body.
\end{proof}

\end{document}
